\section{Design}\label{design}

This chapter explains the strategies used to meet the requirements identified in the analysis. 
%
Ideally, the design should be the same, regardless of the technological choices made during the implementation phase.

\begin{quote}
You can re-arrange the sections as you prefer, but all the sections must be present in the end
\end{quote}

\noindent
\textbf{Important:} try to motivate your design choices in relation to the requirements and features identified 
in \cref{requirements} and \cref{ds-features}.

\subsection{Architecture}\label{architecture}

KompozeR adopts a \textbf{microservice} architectural style behind an
\textbf{API Gateway}, with a \textbf{client--server} interaction between the
browser SPA and the backend. This style was chosen as a direct consequence of
the analysis: the domain decomposes cleanly into loosely-coupled bounded
contexts (\cref{modelling}), and the non-functional requirements of
evolvability and independent scalability (\cref{ds-features}) are best served
by independently deployable services, each owning its own data
(\emph{database-per-service}).

The system is composed of the following logical tiers:

\begin{itemize}
  \item \textbf{Presentation}: a Single Page Application served to the browser.
  It never talks to the services directly; every request goes through the
  gateway.

  \item \textbf{Gateway}: a single entry point that routes REST calls to the
  downstream services, proxies the WebSocket channel, enforces authentication
  and role checks, normalises identity headers, and centralises error handling.
  It also exposes a small number of Backend-for-Frontend (BFF) endpoints that
  aggregate data from several services for specific screens.

  \item \textbf{Domain services}: one service per bounded context ---
  \emph{authentication}, \emph{catalog}, \emph{CAD} (configuration),
  \emph{cart}, \emph{order}, \emph{notification}, \emph{chatbot}, and
  \emph{reporting}. Each service is autonomous, owns its schema, and exposes a
  REST API; a few of them additionally publish or consume asynchronous events.

  \item \textbf{Data and messaging}: a dedicated MongoDB database per service
  and a Redis instance used as an event bus (publish/subscribe) and cache.
\end{itemize}

Three communication styles coexist, each used where it fits best:

\begin{itemize}
  \item \textbf{Synchronous REST} for commands and queries (the default
  interaction);
  \item \textbf{Asynchronous events} over Redis pub/sub for domain events that
  must be propagated to several contexts (notably catalog
  \texttt{PRICE\_CHANGED} and \texttt{AVAILABILITY\_CHANGED});
  \item \textbf{Real-time push} over Socket.IO (WebSocket) for notifications and
  chatbot messages, so users receive updates without polling.
\end{itemize}

Keeping the domain interaction predominantly request/response, and confining
asynchronous and real-time mechanisms to the areas that genuinely benefit from
them (event propagation and push delivery), keeps the distributed complexity
localised and the overall design easier to reason about.

\begin{quote}
Component diagrams are welcome here
\end{quote}

\subsection{Infrastructure}\label{infrastructure}

The system introduces the following infrastructural components:

\begin{itemize}
  \item \textbf{Clients}: web browsers running the SPA.

  \item \textbf{Reverse proxy / static server}: an nginx instance serves the
  static SPA build and reverse-proxies \texttt{/api} (including the WebSocket
  upgrade) to the gateway, acting as the single public entry point of the
  application.

  \item \textbf{API Gateway}: one instance, the only component reachable by
  clients for API traffic.

  \item \textbf{Domain services}: eight stateless backend services
  (authentication, catalog, CAD, cart, order, notification, chatbot,
  reporting).

  \item \textbf{Databases}: one MongoDB deployment per service. Most contexts
  use a single-node MongoDB; the \emph{CAD} context, being the core workflow,
  uses a \textbf{MongoDB replica set} for fault tolerance.

  \item \textbf{Message broker / cache}: a single \textbf{Redis} instance used
  both as a pub/sub event bus between services and as a cache.
\end{itemize}

\paragraph{Distribution over the network.} In the reference deployment all
components run as containers within a single logical environment (a Docker
Compose stack for development, a Kubernetes namespace for a production-like
setup). Clients are remote and reach the system over the Internet; services,
brokers and databases sit on the same cluster network but in separate
containers/pods, so that each can fail, restart, and scale independently. The
CAD replica set is deliberately spread over three separate members (two
data-bearing nodes plus an arbiter in Compose, three data-bearing nodes in
Kubernetes) to tolerate the loss of a node.

\paragraph{Naming and discovery.} Components find each other by logical name
rather than by address. Within the container network, services are reached
through their service name resolved by the platform DNS (Docker Compose's
embedded DNS, or Kubernetes \texttt{Service} DNS), and each service is
configured with the base URLs of its collaborators. Clients only ever know a
single name --- the gateway --- which decouples them from the internal topology
and lets services move or scale without affecting the front end.

\begin{quote}
Component diagrams are welcome here
\end{quote}

\subsection{Modelling}\label{modelling}

\begin{itemize}
  \item which \textbf{domain entities} are there?
  %
  \begin{itemize}
    \item e.g.~\emph{users}, \emph{products}, \emph{orders}, \emph{etc.}
  \end{itemize}
  
  \item how do \emph{domain entities} \textbf{map to} \emph{infrastructural components}?
  %
  \begin{itemize}
    \item e.g.~state of a video game on central server, while inputs/representations on clients
    \item e.g.~where to store messages in an IM app? for how long?
  \end{itemize}

  \item which \textbf{domain events} are there?
  %
  \begin{itemize}
    \item e.g.~\emph{user registered}, \emph{product added to cart}, \emph{order placed}, \emph{etc.}
  \end{itemize}
  
  \item which sorts of \textbf{messages} are exchanged?
  %
  \begin{itemize}
    \item e.g.~\emph{commands}, \emph{events}, \emph{queries}, \emph{etc.}
  \end{itemize}

  \item what information does the \textbf{state} of the system comprehend
  %
  \begin{itemize}
    \item e.g.~\emph{users' data}, \emph{products' data}, \emph{orders' data}, \emph{etc.}
  \end{itemize}
\end{itemize}

\begin{quote}
Class diagram are welcome here
\end{quote}

\subsection{Interaction}\label{interaction}

\begin{itemize}
  \item how do components \emph{communicate}? \emph{when}? \emph{what}?
  \item \emph{which} \textbf{interaction patterns} do they enact?
\end{itemize}

\begin{quote}
Sequence diagrams are welcome here
\end{quote}

\subsection{Behaviour}\label{behaviour}

\begin{itemize}
  \item how does \emph{each} component \textbf{behave} individually (e.g.~in \emph{response} to \emph{events} or messages)?
  \begin{itemize}
    \item some components may be \emph{stateful}, others \emph{stateless}
  \end{itemize}

  \item which components are in charge of updating the \textbf{state} of the system? \emph{when}? \emph{how}?
\end{itemize}

\begin{quote}
State diagrams are welcome here
\end{quote}

\subsection{Data and Consistency Issues}\label{data-and-consistency-issues}

\begin{itemize}
  \item Is there any data that needs to be stored?
  %
  \begin{itemize}
    \item \emph{what} data? \emph{where}? \emph{why}?
  \end{itemize}
  
  \item how should \emph{persistent data} be \textbf{stored}?
  %
  \begin{itemize}
    \item e.g.~relations, documents, key-value, graph, etc.
    \item why?
  \end{itemize}
  
  \item Which components perform queries on the database?
  %
  \begin{itemize}
    \item \emph{when}? \emph{which} queries? \emph{why}?
    \item concurrent read? concurrent write? why?
  \end{itemize}
  
  \item Is there any data that needs to be shared between components?
  %
  \begin{itemize}
    \item \emph{why}? \emph{what} data?
  \end{itemize}
\end{itemize}

\subsection{Fault-Tolerance}\label{fault-tolerance}

\begin{itemize}
  \item Is there any form of data \textbf{replication} / federation / sharing?
  %
  \begin{itemize}
    \item \emph{why}? \emph{how} does it work?
  \end{itemize}
  
  \item Is there any \textbf{heart-beating}, \textbf{timeout}, \textbf{retry mechanism}?
  %
  \begin{itemize}
    \item \emph{why}? \emph{among} which components? \emph{how} does it work?
  \end{itemize}
  
  \item Is there any form of \textbf{error handling}?
  %
  \begin{itemize}
    \item \emph{what} happens when a component fails? \emph{why}? \emph{how}?
  \end{itemize}
\end{itemize}

\subsection{Availability}\label{availability}

\begin{itemize}
  \item Is there any \textbf{caching} mechanism?
  %
  \begin{itemize}
    \item \emph{where}? \emph{why}?
  \end{itemize}
  
  \item Is there any form of \textbf{load balancing}?
  %
  \begin{itemize}
    \item \emph{where}? \emph{why}?
  \end{itemize}

  \item In case of \textbf{network partitioning}, how does the system behave?
  %
  \begin{itemize}
    \item \emph{why}? \emph{how}?
  \end{itemize}
\end{itemize}

\subsection{Security}\label{security}

\begin{itemize}
  \item Is there any form of \textbf{authentication}?
  %
  \begin{itemize}
    \item \emph{where}? \emph{why}?
  \end{itemize}

  \item Is there any form of \textbf{authorization}?
  %
  \begin{itemize}
    \item which sort of \emph{access control}?
    \item which sorts of users / \emph{roles}? which \emph{access rights}?
  \end{itemize}

  \item Are \textbf{cryptographic schemas} being used?
  %
  \begin{itemize}
    \item e.g.~token verification,
    \item e.g.~data encryption, etc.
  \end{itemize}
\end{itemize}
